% ===========================================================================
%  Diagrama a bloques - Practica 3: Medidor de distancia IR (fototransistor)
%  Programa ESP32-S3 (practica3_medidor_distancia.ino)
%  Optoelectronica (V3736) - Dr. Ruben Estrada Marmolejo, CUCEI-UDG
% ===========================================================================
\documentclass[border=12pt]{standalone}
\usepackage[T1]{fontenc}
\usepackage[utf8]{inputenc}
\usepackage[spanish,es-noshorthands]{babel}
\usepackage{amsmath}
\usepackage{tikz}
\usetikzlibrary{positioning, arrows.meta, shapes.geometric, calc, fit, backgrounds}

% ---------- Paleta ----------
\definecolor{colEntrada}{HTML}{1F6FB2}   % entrada / ADC (azul)
\definecolor{colEntradaBg}{HTML}{E3F0FA}
\definecolor{colConv}{HTML}{0E8A6A}      % conversiones (verde)
\definecolor{colConvBg}{HTML}{E1F5EE}
\definecolor{colModelo}{HTML}{B5651D}    % modelos (naranja/ocre)
\definecolor{colModeloBg}{HTML}{FBEEDD}
\definecolor{colSalida}{HTML}{6A3FB5}    % salida (morado)
\definecolor{colSalidaBg}{HTML}{EEE6F8}
\definecolor{colCaptura}{HTML}{B32651}   % captura -> Sheets (magenta)
\definecolor{colCapturaBg}{HTML}{FBE3EC}
\definecolor{colSetup}{HTML}{555F6B}     % setup (gris)
\definecolor{colSetupBg}{HTML}{EDEFF2}
\definecolor{colLinea}{HTML}{2B2F36}

\begin{document}
\begin{tikzpicture}[
  font=\sffamily,
  >=Latex,
  node distance=8mm and 12mm,
  every node/.style={transform shape},
  base/.style={rounded corners=3pt, draw, line width=0.7pt, align=center,
               inner sep=5pt, minimum height=10mm, text width=34mm},
  titulo/.style={font=\sffamily\bfseries},
  entrada/.style={base, draw=colEntrada, fill=colEntradaBg, text=colLinea},
  setup/.style={base, draw=colSetup, fill=colSetupBg, text=colLinea},
  conv/.style={base, draw=colConv, fill=colConvBg, text=colLinea},
  modelo/.style={base, draw=colModelo, fill=colModeloBg, text=colLinea},
  salida/.style={base, draw=colSalida, fill=colSalidaBg, text=colLinea},
  captura/.style={base, draw=colCaptura, fill=colCapturaBg, text=colLinea},
  disp/.style={base, draw=colSetup, fill=white, text=colLinea, text width=30mm},
  flecha/.style={->, line width=0.9pt, draw=colLinea},
  flechaCap/.style={->, line width=0.9pt, draw=colCaptura},
  etiq/.style={font=\sffamily\scriptsize\itshape, text=colLinea, fill=white,
               inner sep=1pt}
]

% ================= FILA SUPERIOR: entrada y setup =================
\node[entrada] (foto) {\textbf{Fototransistor}\\ (colector)\\ {\scriptsize Vout $=$ VCC $-$ Ic$\cdot$Rc}};
\node[entrada, right=of foto] (adc) {\textbf{ADC GPIO10}\\ {\scriptsize ADC1\_CH9}\\ 12 bits, 12 dB};
\node[setup, right=16mm of adc] (setup) {\textbf{setup()}\\ {\scriptsize Serial 115200}\\ {\scriptsize config. ADC + guia}};

% ================= LECTURA =================
\node[conv, below=of adc] (lectura) {\textbf{leerCrudoPromedio()}\\ {\scriptsize promedio de 16}\\ raw (0..4095)};

% ================= CONVERSIONES (dos salidas paralelas del raw) =====
\node[conv, below left=10mm and -2mm of lectura] (vsin)
   {\textbf{V sin calibrar}\\ {\scriptsize $3300\cdot$raw$/4095$}\\ {[mV]}};
\node[conv, below right=10mm and -2mm of lectura] (vcal)
   {\textbf{V calibrado}\\ {\scriptsize analogRead}\\ {\scriptsize MilliVolts()} [mV]};

% ================= LINEALIZACION =================
\node[conv, below=14mm of vcal] (lin)
   {\textbf{Linealizacion}\\ {\scriptsize Vsig $=$ VCC $-$ Vout}\\ {\scriptsize X $=1/\sqrt{\text{Vsig}}$}};

% ================= MODELOS (dos ramas) =================
\node[modelo, below left=14mm and 12mm of lin, text width=32mm] (mlineal)
   {\textbf{Modelo LINEAL}\\ {\scriptsize d $=$ m$\cdot$X $+$ b}\\ {\scriptsize m, b del alumno}};
\node[modelo, below=14mm of lin, text width=32mm] (mcuad)
   {\textbf{Modelo CUADRATICO}\\ {\scriptsize d $=$ aV$^2$+bV+c}\\ {\scriptsize a,b,c de Sheets}};

% ================= SALIDA =================
\node[salida, below=12mm of $(mlineal)!0.5!(mcuad)$] (monitor)
   {\textbf{Monitor Serial}\\ {\scriptsize raw, Vsin, Vcal, Vsig,}\\ {\scriptsize X, d\_lineal, d\_cuad}};

% ================= RAMA DE CAPTURA / CALIBRACION =================
\node[captura, right=44mm of lin, text width=32mm] (captura)
   {\textbf{capturarPunto()}\\ {\scriptsize comando ``punto <cm>''}\\ {\scriptsize fila CSV:}\\ {\scriptsize d\_real, raw, V\_cal, X}};
\node[captura, below=9mm of captura, text width=32mm] (sheets)
   {\textbf{Google Sheets}\\ {\scriptsize minimos cuadrados}\\ {\scriptsize $\to$ m, b y cuadratica}};
\node[captura, below=9mm of sheets, text width=32mm] (escribe)
   {\textbf{Se escriben coef.}\\ {\scriptsize en el programa}\\ {\scriptsize (PENDIENTE/OFFSET, coefA..C)}};

% ================= DESPACHADOR (loop) =================
\node[disp, above=13mm of captura, text width=38mm] (loop)
   {\textbf{loop(): despachador}\\ {\scriptsize lee Serial}\\ {\scriptsize leer / stream / punto <cm> / coef / ayuda}};

% ===================== FLECHAS DE FLUJO =====================
\draw[flecha] (foto) -- (adc);
\draw[flecha] (adc) -- (lectura);

% raw -> dos conversiones
\draw[flecha] (lectura.south) -- ++(0,-3mm) -| (vsin.north);
\draw[flecha] (lectura.south) -- ++(0,-3mm) -| (vcal.north);

% V calibrado -> linealizacion
\draw[flecha] (vcal) -- (lin);

% linealizacion -> modelo lineal (usa X)
\draw[flecha] (lin.south) -- ++(0,-3mm) -| (mlineal.north)
     node[etiq, pos=0.25] {X};
% V calibrado -> modelo cuadratico (usa V, no X)
\draw[flecha] (vcal.east) -- ++(0.8,0) |- (mcuad.east)
     node[etiq, pos=0.25] {V};

% modelos -> monitor
\draw[flecha] (mlineal.south) |- ($(monitor.north)+(-8mm,3mm)$) -- ($(monitor.north)+(-8mm,0)$);
\draw[flecha] (mcuad.south)   |- ($(monitor.north)+( 8mm,3mm)$) -- ($(monitor.north)+( 8mm,0)$);

% V sin calibrar -> monitor (se reporta directo)
\draw[flecha] (vsin.west) .. controls +(-1.4,0) and +(-1.4,0) .. (monitor.west)
     node[etiq, pos=0.5] {se reporta};

% ---- rama de captura ----
\draw[flechaCap] (loop) -- (captura);
\draw[flechaCap] (lectura.east) .. controls +(2.0,0) and +(0,1.4) .. (captura.north)
     node[etiq, pos=0.5] {raw};
\draw[flechaCap] (captura) -- (sheets) node[etiq, pos=0.5] {pegar CSV};
\draw[flechaCap] (sheets) -- (escribe);
% coeficientes vuelven a los modelos
\draw[flechaCap, dashed] (escribe.south) -- ++(0,-4mm) -| ([xshift=6mm]mcuad.south)
     node[etiq, pos=0.25] {coef. finales};

% ===================== TITULO Y LEYENDA =====================
\node[titulo, font=\sffamily\bfseries\large, text=colLinea, above=2mm of foto.north west, anchor=south west]
   (tit) {Programa ESP32-S3 -- Medidor de distancia IR (Practica 3, Optoelectronica)};

% Fondo agrupador de la rama de captura/calibracion
\begin{scope}[on background layer]
  \node[rounded corners=4pt, fill=colCapturaBg!45, draw=colCaptura!40, dashed,
        inner sep=6mm, fit=(loop)(captura)(sheets)(escribe),
        label={[colCaptura,font=\sffamily\bfseries\small]above:CALIBRACION (menu + Google Sheets)}] {};
\end{scope}

\end{tikzpicture}
\end{document}
